
\chapter{Sessió 4 (05/10/2026). Els àtoms existeixen. Solucions.}

Aquestes són les fórmules de les dues substàncies que hem utilitzat:

\[
\text{Bicarbonat de sodi:}\qquad \mathrm{NaHCO_3}
\]

\[
\text{Àcid cítric:}\qquad \mathrm{C_6H_8O_7}
\]

El bicarbonat de sodi conté, per cada unitat de la substància:

\begin{itemize}
    \item 1 àtom de sodi, \(\mathrm{Na}\);
    \item 1 àtom d'hidrogen, \(\mathrm{H}\);
    \item 1 àtom de carboni, \(\mathrm{C}\);
    \item 3 àtoms d'oxigen, \(\mathrm{O}\).
\end{itemize}

El nom \emph{bicarbonat} és una mica traïdor: aquest \emph{bi-} no vol dir que hi hagi dos àtoms de carboni. És un nom que ve de la nomenclatura química antiga. La fórmula és qui mana, i aquí només hi ha un carboni.

I això de \(\mathrm{Na}\) per al sodi també té història. No sembla tenir gaire relació amb ``sodi'', oi? Pots investigar d'on surt aquest símbol.

L'àcid cítric conté:

\begin{itemize}
    \item 6 àtoms de carboni, \(\mathrm{C}\);
    \item 8 àtoms d'hidrogen, \(\mathrm{H}\);
    \item 7 àtoms d'oxigen, \(\mathrm{O}\).
\end{itemize}

Es troba de manera natural en moltes fruites i també es pot fabricar industrialment. Juntament amb el sucre és, a més, una de les bases d'aquelles espantoses llaminadures àcides que semblen dissenyades per comprovar fins on pot arribar l'esmalt dental.

\section{La taula periòdica serveix per a alguna cosa}

Recordem ara la taula periòdica.

No serveix només per tenir un pòster científic penjat a la paret.

Els símbols de les fórmules ens indiquen amb quins àtoms està construïda cada substància, una mica com les instruccions d'un LEGO. I la taula periòdica ens dona, entre moltes altres coses, la massa de cadascun d'aquests àtoms.

Per tant, podem sumar-los.

Per al bicarbonat de sodi:

\[
\mathrm{NaHCO_3}
\]

\[
22,99 + 1,01 + 12,01 + 3(16,00)
\simeq 84,01
\]

La seva massa molar és, per tant:

\[
84,01\ \text{g/mol}
\]

Per a l'àcid cítric:

\[
\mathrm{C_6H_8O_7}
\]

\[
6(12,01)+8(1,01)+7(16,00)
\simeq 192,12
\]

i obtenim:

\[
192,12\ \text{g/mol}
\]

Ja tenim una manera de connectar aquelles fórmules diminutes amb els grams que acabem de pesar a la balança.

\section{Ara fem-los reaccionar}

En dissoldre les substàncies en aigua permetem que les partícules es moguin i entrin en contacte amb molta més facilitat, de manera que la reacció es produeix ràpidament.

Sense entrar encara en tots els detalls, el que passa es pot escriure així:

\[
\mathrm{C_6H_8O_7
+3\,NaHCO_3
\longrightarrow
Na_3C_6H_5O_7
+3\,CO_2
+3\,H_2O}
\]

Aquesta és la primera vegada que veiem una \textbf{equació química}.

A partir d'ara les trobarem fins a la sopa.

L'equació ens diu moltes coses, però de moment fixa't sobretot en això:

\[
\mathrm{3\,CO_2}
\]

Aquest \(\mathrm{CO_2}\) és el diòxid de carboni.

Sí: són les bombolletes.

I les bombolletes han marxat del recipient i se n'han anat a l'aire.

\section{Fem els números}

L'equació ens diu que una molècula d'àcid cítric reacciona amb tres unitats de bicarbonat de sodi.

En masses molars:

\[
192,12\ \text{g d'àcid cítric}
\]

reaccionen amb:

\[
3\times84,01=252,03\ \text{g de bicarbonat}
\]

i produeixen:

\[
3\times44,01=132,03\ \text{g de CO}_{2}
\]

Nosaltres no hem utilitzat \(192,12\ \text{g}\) d'àcid cítric.

N'hem utilitzat només:

\[
10,0\ \text{g}
\]

Fem una regla de tres:

\[
192,12\ \text{g d'àcid cítric}
\longrightarrow
132,03\ \text{g de CO}_{2}
\]

\[
10,0\ \text{g d'àcid cítric}
\longrightarrow
x
\]

i obtenim:

\[
x=
\frac{10,0\times132,03}{192,12}
\simeq 6,87\ \text{g}
\]

O sigui:

\[
\boxed{6,9\ \text{g de CO}_{2}}
\]

Just el que acabem de veure desaparèixer de la balança.

I també podem calcular quant bicarbonat necessitem exactament per reaccionar amb aquests \(10\ \text{g}\) d'àcid cítric:

\[
\frac{10,0\times252,03}{192,12}
\simeq 13,1\ \text{g}
\]

Per això un grup tenia exactament \(13,1\ \text{g}\).

Si n'hi posem \(15\ \text{g}\), simplement sobra bicarbonat. I posar \(50\) o \(60\ \text{g}\) d'aigua tampoc no crea més diòxid de carboni: l'aigua ajuda a fer possible la reacció, però no canvia quants àtoms d'àcid cítric tenim disponibles.

Quan l'àcid cítric s'acaba, s'ha acabat la festa.

\section{Els àtoms existeixen}

Ja ho tenim.

Si acceptem que els àtoms existeixen i que les substàncies estan construïdes amb nombres concrets d'aquests àtoms, el resultat de l'experiment encaixa gairebé peça per peça.

Hem començat amb \(10\ \text{g}\) d'àcid cítric.

La teoria prediu que podem produir aproximadament:

\[
6,9\ \text{g de CO}_{2}
\]

I la balança ens diu que han desaparegut aproximadament:

\[
6,9\ \text{g}
\]

Hauríem de tirar bastant d'imaginació per trobar una altra explicació que encaixés igual de bé.

En certa manera estem repetint el camí que van seguir científics com Dalton, Antoine Lavoisier i Marie-Anne Pierrette Paulze Lavoisier: pesar substàncies amb molta cura, observar en quines proporcions reaccionen i intentar entendre què ens estan dient aquests números.

Estem connectant amb els primers passos de la química moderna i repetint-los nosaltres mateixos.

Només que aquesta vegada hem començat amb un volcà de cuina.